%-------------------------------------------------------------------------------
% Colter Snyder -- Cover letter, Google Solutions Engineer, Chrome Scaled Solutions,
% gUP Chrome and Web Ecosystem, gTech
% Build: latexmk -pdf -outdir=build ColterSnyder_CoverLetter_Google_ChromeSE.tex
%-------------------------------------------------------------------------------
\documentclass[letterpaper,11pt]{article}

\usepackage[margin=0.8in,top=0.6in,bottom=0.6in]{geometry}
\usepackage[T1]{fontenc}
\usepackage{lmodern}
\usepackage{xcolor}
\usepackage{fontawesome5}
\usepackage[hidelinks]{hyperref}

\definecolor{accent}{HTML}{1F5FA8}
\pagestyle{empty}
\setlength{\parindent}{0pt}
\setlength{\parskip}{8pt}
\pdfgentounicode=1

\begin{document}

%--- Header (matches resume) ---------------------------------------------------
\begin{center}
  {\Huge\bfseries Colter Snyder}\par\vspace{4pt}
  \footnotesize
  Lakewood, CO \enspace$\cdot$\enspace
  \faPhone*\ (720) 626-2558 \enspace$\cdot$\enspace
  \faEnvelope\ \href{mailto:coltersnyder@gmail.com}{coltersnyder@gmail.com} \enspace$\cdot$\enspace
  \mbox{\faLinkedin\ \href{https://www.linkedin.com/in/colter-snyder-a36854183/}{in/colter-snyder-a36854183}}
\end{center}
{\color{accent}\hrule height 0.6pt}
\vspace{6pt}

October 23, 2026

Hiring Team\\
Google

\textbf{Re: Solutions Engineer, Chrome Scaled Solutions, gUP Chrome and Web Ecosystem}

Dear Hiring Team,

The web is moving from pages people browse to services that agents act on, and Chrome sits at the
center of that change. I would like to help partners and developers make that shift. I am a software
engineer at Innovative Scientific Solutions, where I build web applications for defense customers, and
I hold an M.S.\ in Computer Science from the Colorado School of Mines.

I have built for the web since 2021, starting as the webmaster for Mines' computer science outreach
site in HTML, CSS, and JavaScript. Today I develop Angular and TypeScript interfaces, contribute heavily
to UI/UX design, and build the Python and Rust services behind them. I also spearheaded our internal
performance-benchmarking tool, a Python command-line application with a GUI front end that runs any
combination of our 13 effect calculations and charts the results. It exposed a tenfold slowdown in how
we call an external library, and its graphs made the problem easy to present to leadership. That project taught me how much
faster a team moves when the data is in front of it, which is the kind of feedback loop this role
builds between partners, users, and Chrome's product teams.

AI-assisted development is part of how I work. I use Claude Code and Cursor regularly, and because my
current work is export-controlled, I have had to think carefully about where these tools fit: I use
them only for non-controlled exploration and research, never with controlled code or data. Evaluating
AI-assisted onboarding and troubleshooting workflows calls for the same judgment about where
automation helps and where it needs guardrails.

The role also asks for teaching. I co-taught Data Structures and Algorithms to 120 students, taught
another section of 90, designed STEM summer camps, and co-presented a teacher workshop, so turning
technical knowledge into playbooks and guidance that others can act on is work I enjoy. I would welcome
the chance to bring that to the Chrome Scaled Solutions team. I live in Lakewood and would be glad to work from Google's Boulder office or remotely. Thank you for your time and consideration.

Sincerely,

Colter Snyder

\end{document}
